% Created 2026-10-21 Wed 19:48
% Intended LaTeX compiler: pdflatex
\documentclass[11pt]{article}
\usepackage[utf8]{inputenc}
\usepackage[T1]{fontenc}
\usepackage{graphicx}
\usepackage{longtable}
\usepackage{wrapfig}
\usepackage{rotating}
\usepackage[normalem]{ulem}
\usepackage{amsmath}
\usepackage{amssymb}
\usepackage{capt-of}
\usepackage{hyperref}
\documentclass[10pt,a4paper,showtrims]
\hypersetup{colorlinks=true, urlcolor={blue}, linkcolor={blue}}
\usepackage[labelfont=bf]{caption}
\usepackage{tcolorbox}
\author{Triscuitcircuit}
\date{\today}
\title{macOS OpenXR game example}
\hypersetup{
 pdfauthor={Triscuitcircuit},
 pdftitle={macOS OpenXR game example},
 pdfkeywords={},
 pdfsubject={},
 pdfcreator={Emacs 30.2 (Org mode 9.7.11)}, 
 pdflang={English}}
\begin{document}

\maketitle
\tableofcontents

\clearpage

A small target game: aim a tracked controller at the sphere and squeeze the trigger to score. The camera and targets use world coordinates; the game does not move the world with head rotation. There is no artificial locomotion.
\section{Requirements}
\label{sec:org6191da0}
\begin{itemize}
\item Godot 4 with the OpenXR module and Vulkan renderer available on macOS. An engine built without XR will not work; if building from source, enable \texttt{module\_openxr\_enabled=yes} and \texttt{disable\_xr=no}.
\item A compatible, separately supplied macOS OpenXR runtime providing Vulkan graphics, tracked headset/controllers, and a working USB headset transport.
\item Matching architectures for the engine, runtime, Vulkan implementation, and all dynamic libraries (for example all arm64).
\item A connected and awake headset. Close any dashboard/application holding an exclusive runtime session before starting the game.
\end{itemize}

This example uses the engine's default OpenXR action map, including \texttt{trigger}, \texttt{left\_hand}, \texttt{right\_hand}, and aim poses. If you customize it, retain these actions/bindings or update the script.
\section{Run}
\label{sec:org058a7e5}
\begin{verbatim}
python3 launch.py --runtime /absolute/path/to/runtime.json --godot /path/to/Godot.app/Contents/MacOS/Godot
\end{verbatim}
Add \texttt{-{}-{}icd /absolute/path/to/MoltenVK\_icd.json} if your runtime deployment requires it. For an exported macOS game, use \texttt{-{}-{}game /path/to/Game.app/Contents/MacOS/Game} instead of \texttt{-{}-{}godot}.

The launcher uses argument arrays, validates the manifest, and sets the runtime before starting the engine. OpenXR and XR shaders are enabled in project settings. Use Mobile + Vulkan for this runtime path rather than the native Metal renderer.
\section{How Newt fits}
\label{sec:org878d55f}
Godot owns its OpenXR session and render loop. \textbf{Do not call \texttt{Newt\_Run()} inside this Godot game:} Newt's current dashboard also owns a session and rendering loop. Use the dashboard as a separate launcher and hand off the headset to Godot. This example uses the external runtime directly; it does not embed Newt's dashboard.

\texttt{Newt::usb} is available to native host tools for generic USB operations. Do not claim the headset's USB interface from a game while its runtime owns it. libusb transport alone does not implement Quest pairing, tracked poses, encoded video or timewarp.

Embedding Newt panels directly in Godot would need a new GDExtension that exposes capture frames and panel state as Godot textures/nodes, uses Godot's existing OpenXR session, and avoids starting Newt's own Vulkan renderer. That adapter is not implemented here.
\section{Bundle an exported macOS game}
\label{sec:orge99d38e}
A deployment layout can be:
\begin{verbatim}
Game.app/Contents/
  MacOS/Game                         # exported Godot executable
  Resources/Game.pck                 # as generated by the export
  Resources/OpenXR/runtime.json
  Resources/OpenXR/MoltenVK_icd.json  # only if using an external ICD
  Frameworks/XR/                     # external runtime and its dependency dylibs
  Resources/Licenses/                # notices for every shipped component
\end{verbatim}
\begin{enumerate}
\item Export the project using matching macOS export templates. Preserve the export's actual executable and resource filenames.
\item Supply your runtime as a separately built artifact with its dependency list, license notices, source-distribution materials where required, and supported headset/architecture documented. This repository does not contain that artifact.
\item Copy its dylibs into \texttt{Frameworks/XR}. Rewrite non-system load paths to bundle-relative \texttt{@loader\_path} / \texttt{@rpath} paths. The runtime JSON's \texttt{runtime.library\_path} can point to \texttt{../../Frameworks/XR/libopenxr\_monado.dylib} if that is the selected runtime's real filename. Rewrite an external ICD manifest to its actual bundled library path too.
\item Use a small native launcher as the app's \texttt{CFBundleExecutable} to resolve the bundle path, set \texttt{XR\_RUNTIME\_JSON} (and the ICD selection when needed), then exec the exported game executable. \texttt{launch.py} demonstrates the environment/argument contract for development; do not depend on users having Python installed for a shipping app. Finder launches do not inherit variables from your terminal.
\item Sign nested libraries and executables, then the outer app; notarize your release as appropriate. Verify all dependencies with \texttt{otool -L}, and test on a clean Mac without Homebrew or developer environment variables.
\end{enumerate}

Do not assume changing a build flag removes GPL code: audit the actual runtime artifact. Newt's MIT license does not relicense runtime code or libusb.
\section{Validation status}
\label{sec:org2f8a168}
The project can be imported and run without XR for script/scene checks. A successful desktop check does not validate Vulkan/OpenXR initialization, controller bindings, frame timing, USB transport, or headset audio. Those need a hardware run with the selected engine/runtime combination.
\section{References}
\label{sec:org1bff404}
\begin{itemize}
\item \href{https://docs.godotengine.org/en/stable/tutorials/xr/setting\_up\_xr.html}{Godot XR setup}
\item \href{https://github.com/godotengine/godot/blob/master/modules/openxr/config.py}{Godot OpenXR build configuration}
\end{itemize}
\section{Editor Play integration}
\label{sec:org271ada3}
Use the Newt Play Bridge plugin to discover/start Newt automatically on Play, hand off the headset, and return after the game exits. The session manager also creates development app bundles that reuse the installed runtime.
\end{document}
